% Synthetic fixture: a short single-file article.
\documentclass[11pt]{article}
\usepackage[utf8]{inputenc}
\usepackage{amsmath}
\newcommand{\method}{\textsc{Weave}}
\title{A Small Study of Merge Latency}
\author{A. Researcher \and B. Colleague}
\date{}

\begin{document}
\maketitle

\begin{abstract}
We measure how long concurrent edits take to converge in a shared document.
Latency stays below 40~ms for 95\% of edits, and \emph{no acknowledged edit is lost}.
\end{abstract}

\section{Introduction}
\label{sec:intro}

Collaborative editors must merge concurrent changes without losing work~\cite{shapiro2011}.
Earlier systems used a central transform server; later ones use replicated data
types \cite{nicolaescu2016,shapiro2011}. % TODO: add a third citation
We call our prototype \method{} and describe it in Section~\ref{sec:method}.

Our contributions are:
\begin{itemize}
  \item a merge protocol that acknowledges an edit only after it is durable;
  \item a measurement of latency under \textbf{three} concurrent writers;
  \item an open fixture set, described in Section~\ref{sec:results}.
\end{itemize}

\section{Method}
\label{sec:method}

Each replica holds a sequence $s = (c_1, \dots, c_n)$ of characters. An edit is a
pair $(p, \delta)$ where $p$ is a position and $\delta$ the inserted text.
The expected latency is
\begin{equation}
  \label{eq:latency}
  L = t_{\text{persist}} + t_{\text{ack}} + t_{\text{broadcast}} .
\end{equation}
Equation~\eqref{eq:latency} ignores client rendering, which we measure separately.

\subsection{Setup}

We ran every trial on one machine -- no network hops -- with a 2~GHz budget
per replica. The ``slow path'' covers retries \& reconnects.

\section{Results}
\label{sec:results}

Table~\ref{tab:latency} reports the median and tail latency.

\begin{table}[h]
  \centering
  \begin{tabular}{lrr}
    Writers & Median (ms) & p95 (ms) \\
    \hline
    1 & 4 & 9 \\
    3 & 11 & 38 \\
  \end{tabular}
  \caption{Merge latency by number of concurrent writers.}
  \label{tab:latency}
\end{table}

Latency grows with the number of writers but stays under our 40~ms target.

\bibliographystyle{plain}
\bibliography{refs}

\end{document}
